\subsection{A central involution forced by two twist moments}
\label{sec:unbounded-fixed-orbits}
Consider a self-dual rank-four ring with exactly two invertible basis
objects $1,g$, where $g$ fixes the other two objects $X,Y$. Its products are
\begin{equation}\label{eq:fixed-orbit-rules}
\begin{gathered}
 g^2=1,\qquad gX=X,\qquad gY=Y,\\
 X^2=1+g+aX+bY,\qquad XY=bX+cY,\\
 Y^2=1+g+cX+eY,
 \qquad 2+ac+be=b^2+c^2.
\end{gathered}
\end{equation}
All parameters are nonnegative integers.

\begin{lemma}\label{lem:fixed-orbits-galois}
A ring of the form \eqref{eq:fixed-orbit-rules} admits a pivotal
categorification if and only if it admits a pseudo-unitary categorification.
\end{lemma}
\begin{proof}
The rational character $(1,-1,0,0)$ splits off a copy of $\mathbb Q$ from
$R\otimes\mathbb Q$. Its complementary algebra has dimension three.
Consequently the degree $k$ of the Frobenius--Perron character field is
$1$, $2$, or $3$. A pivotal categorification is spherical because every
simple is self-dual; its dimension character cannot be $(1,-1,0,0)$.

If $k=1$, its Frobenius--Perron dimension is an integer, and the category is
pseudo-unitary. If $k=3$, all three remaining characters form one Galois
orbit, so \cite[Lemma~\ref{B-lem:galois}]{Finite} applies. Suppose $k=2$ and write
$D=\operatorname{FPdim}(R)$. If $D$ is rational, it is an integer and the
same weak-integrality argument applies. Otherwise $D$ generates the
quadratic character field. Besides the two conjugate Frobenius--Perron
characters, the only remaining character is rational. It cannot be the
spherical dimension character, since the spherical dimension field must
contain every formal codegree, including the nonrational number $D$.
Thus the dimension character is Galois conjugate to Frobenius--Perron,
which proves the assertion.
\end{proof}

\begin{proposition}\label{prop:fixed-orbits-boson}
Suppose $b,c>0$ and that \eqref{eq:fixed-orbit-rules} admits a
pseudo-unitary categorification $\mathcal C$. Then the object $g$ lifts
to an invertible object of order two and twist one in $\mathcal Z(\mathcal C)$.
Moreover $\mathcal C$ is a $C_2$-equivariantization of a rank-five fusion
category on which the nontrivial group element fixes only the tensor unit.
\end{proposition}
\begin{proof}
Use the canonical spherical structure and put $D=\operatorname{FPdim}(\mathcal C)$.
The character $(1,-1,0,0)$ has formal codegree $2$. Thus
\[
 I(1)=1\oplus A\oplus B\oplus C,
 \qquad \dim(B)=D/2,
 \qquad \dim(A)+\dim(C)=D/2-1.
\]
All four summands have twist one. The coefficient of $g$ in $FI(1)$ is
two, while $\dim\operatorname{End}(I(g))=4$, as follows directly from
$FI(g)=\sum_U U\otimes g\otimes U^*$.

The first twist moment gives that the dimension of the low part of $I(g)$
is at most $D/2$. If $g$ occurred in $F(B)$, it would already contribute
$D/2$, and the other required occurrence of $g$ would make this bound
strictly fail. Hence $[F(B):g]=0$. The two occurrences of $g$ are therefore
in $F(A)$ and $F(C)$. They cannot both be in one column: that multiplicity
two would exhaust $\dim\operatorname{End}(I(g))=4$, so $I(g)$ would have
only twist-one constituents, contrary to its vanishing first twist trace.
Thus both multiplicities are one. The low dimension is
\[
 L=D/2-1.
\]
The residual value of $\dim\operatorname{End}(I(g))$ is two, so there
are exactly two further simple summands, each with multiplicity one.
Denote their positive dimensions by $\alpha,\beta$ and twists by $u,v$.
The zeroth, first and second twist moments give
\[
 \alpha+\beta=D/2+1,\qquad
 \alpha u+\beta v=-(D/2-1),\qquad
 \alpha u^2+\beta v^2=D\nu_2(g)-(D/2-1).
\]
If $\nu_2(g)=-1$, the last right-hand side has absolute value
$3D/2-1>D/2+1$, since $D>2$, violating the triangle inequality.
Therefore $\nu_2(g)=1$. Equality in the triangle inequality for the last
identity forces $u^2=v^2=1$. The first identity involving twists then
forces, after interchanging the two summands,
\[
 (\alpha,u)=(1,1),\qquad (\beta,v)=(D/2,-1).
\]
Let $z$ be the central simple of dimension one. Its restriction contains
$g$, and consequently $F(z)=g$.

There is no nontrivial central invertible whose underlying object is the
unit. Such an object would give a tensor automorphism of the identity of
$\mathcal C$, with scalars $\lambda_g,\lambda_X,\lambda_Y$.
The relations $gX=X$, $Y\subset X^2$, and $X\subset Y^2$, together with
self-duality, force respectively
$\lambda_g=1$, $\lambda_Y=\lambda_X^2=1$, and
$\lambda_X=\lambda_Y^2=1$. Hence $z^2=1$.
The order-two invertible $z$ has dimension and twist one, so it generates
a Tannakian subcategory $\operatorname{Rep}(C_2)$ of the center, embedded
fully faithfully into $\mathcal C$ by the forgetful functor.

De-equivariantization gives a fusion category $\mathcal D$ with a
$C_2$-action and $\mathcal C\simeq\mathcal D^{C_2}$.
For clarity, the orbit description of equivariantization determines its
simple objects without a dimension estimate. A fixed simple of
$\mathcal D$ gives two simples of $\mathcal C$, interchanged by tensoring
with the sign representation $g$, because $H^2(C_2,\mathbb C^\times)=0$.
A two-element orbit gives one simple fixed by $g$. In \(\mathcal C\), the
unit and $g$ form the sole two-element orbit under tensoring with $g$;
$X$ and $Y$ are fixed. Thus the action on $\operatorname{Irr}(\mathcal D)$
has exactly one fixed simple, the unit, and two orbits of size two.
In particular $\mathcal D$ has rank five.
\end{proof}

\begin{theorem}\label{thm:fixed-orbits-unbounded}
Among rings of the form \eqref{eq:fixed-orbit-rules}, the only one admitting
a pivotal categorification is $K_0(\operatorname{Rep}(D_{10}))$.
There is no multiplicity restriction.
\end{theorem}
\begin{proof}
By Lemma~\ref{lem:fixed-orbits-galois} we can assume pseudo-unitarity.
First suppose $b,c>0$ and form $\mathcal D$ as in
Proposition~\ref{prop:fixed-orbits-boson}. The dimensions of its simples are
\[
 1,u,u,v,v,\qquad
 u=\operatorname{FPdim}(X)/2,\quad
 v=\operatorname{FPdim}(Y)/2.
\]
They are integers, by the fixed-point-free automorphism theorem
\cite[Theorem~2.3]{SchopierayFPF}.
There is also a short rank-five verification of this integrality.
The complexified fusion algebra is either $\mathbb C^5$ or
$\mathbb C\oplus M_2(\mathbb C)$. In the second case its sole
one-dimensional character, Frobenius--Perron dimension, is Galois fixed.
In the first case, the involution has trace one in the based basis and
therefore fixes exactly one primitive idempotent, hence exactly one
character. Every Galois conjugate of Frobenius--Perron is fixed by the
involution, so it must be that same character. Its values are rational
algebraic integers in either case.

An integral fusion ring with at most three distinct basis dimensions has
a nontrivial pointed subring \cite[Lemma~4.2]{SchopierayFPF}.
The pointed part of $K_0(\mathcal D)$ is invariant under the involution,
so it has rank three or five. If its rank is three, the two remaining
objects have a common integer dimension $v>1$ and are fixed by the
three invertibles. Their product with their dual therefore gives
$v^2=3+kv$ for an integer $k\geq0$, forcing $v=3$. After exchanging $X,Y$ if necessary, their dimensions in
$\mathcal C$ would then be $2,6$. The equation for $X^2$ would read
$4=2+2a+6b$, contradicting $b>0$. If its pointed rank is five, the dimensions of
$\mathcal C$ are $(1,1,2,2)$. The dimension equations in
\eqref{eq:fixed-orbit-rules} force
$(a,b,c,e)=(0,1,1,0)$, the representation ring of $D_{10}$.

It remains to treat a zero among $b,c$. By exchanging $X,Y$, assume
$b=0$. Associativity reads $c(c-a)=2$, so $c=2$, $a=1$. Thus the ring
is an extension of $K_0(\operatorname{Rep}(S_3))$ from
Section~\ref{sec:near-s3}, with arbitrary coefficient $\kappa=e\geq0$.
Put $d=(\kappa+\sqrt{\kappa^2+24})/2$ and $D=12+\kappa d$.
For $\kappa\notin\{0,1,5\}$, this $D$ is quadratic, with minimal polynomial
\[
 t^2-(\kappa^2+24)t+6(\kappa^2+24).
\]
The $d$-number criterion implies $6\mid\kappa^2$, hence $6\mid\kappa$;
these parameters, including $\kappa=0$, are excluded by
Section~\ref{sec:near-s3}. The only positive $\kappa$ for which
$\kappa^2+24$ is a square are $1$ and $5$, as follows by factoring
$(s-\kappa)(s+\kappa)=24$. For the two remaining parameters $\kappa\in\{1,5\}$, the
bounded-census verifier matches the complete tensors and checks their
nonexistence certificates; see \cite[Theorem~\ref{B-thm:finite}]{Finite}.
The finite theorem is independent of the present unrestricted argument.
This completes the proof.
\end{proof}

\begin{corollary}\label{cor:nontrivial-invertible-unbounded}
A rank-four fusion ring with a nontrivial invertible basis element admits
a pivotal categorification if and only if it is one of the eleven rings
in Theorem~\ref{thm:two-dimensions-unbounded}. These all admit unitary
categorifications.
\end{corollary}
\begin{proof}
For a non-self-dual ring with a nontrivial invertible, its noninvertibles
are either a single object or a dual pair, and therefore have equal
Frobenius--Perron dimensions. Theorem~\ref{thm:two-dimensions-unbounded}
applies. In the self-dual case, four invertibles give a pointed ring;
three invertibles cannot occur because a group of order three is not
self-dual. With two invertibles, their action either interchanges the two
remaining objects, which have equal dimensions and are covered by
Theorem~\ref{thm:two-dimensions-unbounded}, or fixes both, which is exactly
Theorem~\ref{thm:fixed-orbits-unbounded}.
\end{proof}
